
\documentclass[12pt,a4paper]{article}

% ===== CÁC GÓI CẦN THIẾT =====
\usepackage[a4paper,margin=2.5cm]{geometry}
\usepackage{fontspec}
\usepackage{polyglossia}
\setdefaultlanguage{vietnamese}
\setmainfont{Times New Roman}

\usepackage{amsmath,amssymb,bm}
\usepackage{enumitem}
\usepackage{setspace}

\setstretch{1.25}
\setlength{\parindent}{1cm}
\setlength{\parskip}{4pt}

\begin{document}

% Đánh số bắt đầu từ mục 3
\setcounter{section}{2}

\section{Gaussian Mixture Model}

\subsection{Model Formulation}

\subsubsection*{Mixture of K Gaussians}

Gaussian Mixture Models (GMM) là một thuật toán phân cụm,
tính phân phối xác suất của mỗi điểm thuộc về tất cả các cụm.
Khác với K-means là mỗi điểm chỉ thuộc về một cụm, GMM
ước lượng xác suất một điểm dữ liệu thuộc về từng cụm.
Điều này cho phép mô hình biểu diễn các cụm có vùng chồng lấn
và thể hiện sự không chắc chắn trong việc phân cụm.

GMM giả định rằng dữ liệu được sinh ra từ hỗn hợp gồm
$K$ phân phối Gaussian đa chiều. Hàm mật độ xác suất của
điểm dữ liệu $\bm{x}_n$ được xác định theo công thức sau:

\begin{equation}
p(\bm{x}_n \mid \boldsymbol{\theta})
=
\sum_{k=1}^{K}
\pi_k
\mathcal{N}(\bm{x}_n \mid \bm{\mu}_k,\boldsymbol{\Sigma}_k)
\end{equation}

Trong đó:
\begin{itemize}[leftmargin=1.5cm]
    \item $K$ là số thành phần Gaussian.
    \item $\pi_k$ là trọng số hỗn hợp của thành phần thứ $k$.
    \item $\bm{\mu}_k$ là vector trung bình.
    \item $\boldsymbol{\Sigma}_k$ là ma trận hiệp phương sai.
    \item $\boldsymbol{\theta}$ đại diện cho toàn bộ tham số
    của mô hình.
\end{itemize}

Các trọng số hỗn hợp thỏa mãn điều kiện:

\begin{equation}
\sum_{k=1}^{K}\pi_k=1,
\qquad \pi_k \geq 0
\end{equation}

\subsubsection*{Multivariate Gaussian Distribution}

Các điểm dữ liệu thường có nhiều đặc trưng, nên GMM
thường sử dụng phân phối Gaussian đa chiều. Công thức
hàm mật độ xác suất của phân phối Gaussian đa chiều được
xác định như sau:

\begin{equation}
\mathcal{N}(\bm{x}\mid\bm{\mu},\boldsymbol{\Sigma})
=
\frac{1}
{(2\pi)^{D/2}|\boldsymbol{\Sigma}|^{1/2}}
\exp\left(
-\frac{1}{2}
(\bm{x}-\bm{\mu})^{T}
\boldsymbol{\Sigma}^{-1}
(\bm{x}-\bm{\mu})
\right)
\end{equation}

Trong đó:
\begin{itemize}[leftmargin=1.5cm]
    \item $D$ là số chiều dữ liệu.
    \item $|\boldsymbol{\Sigma}|$ là định thức của ma trận
    hiệp phương sai.
    \item $\bm{\mu}$ là vector trung bình.
    \item $\boldsymbol{\Sigma}^{-1}$ là ma trận nghịch đảo
    của ma trận hiệp phương sai.
\end{itemize}

$\bm{\mu}$ xác định tâm của một thành phần Gaussian,
trong khi $\boldsymbol{\Sigma}$ mô tả mức độ phân tán
và mối quan hệ giữa các đặc trưng. Nhờ đó, GMM có thể
biểu diễn các cụm có kích thước, hình dạng và hướng
phân bố khác nhau.

\subsubsection*{Parameters: $\pi$, $\mu$, $\Sigma$}

\paragraph{Mixture weights ($\pi_k$)}

Trọng số hỗn hợp biểu thị xác suất tiên nghiệm của việc
lựa chọn thành phần Gaussian thứ $k$. Trọng số càng lớn
thì thành phần tương ứng đóng góp càng nhiều vào phân phối
hỗn hợp tổng thể.

\paragraph{Mean vectors ($\bm{\mu}_k$)}

Vector trung bình xác định vị trí tâm của một thành phần
Gaussian trong không gian đặc trưng. Mỗi thành phần có
một vector trung bình riêng, được ước lượng từ dữ liệu
quan sát trong quá trình huấn luyện mô hình.

\paragraph{Covariance matrices ($\boldsymbol{\Sigma}_k$)}

Ma trận hiệp phương sai mô tả mức độ phân tán và mối
tương quan giữa các đặc trưng trong một thành phần.
Các phần tử trên đường chéo chính biểu thị phương sai
của từng đặc trưng, trong khi các phần tử ngoài đường
chéo biểu thị hiệp phương sai giữa các đặc trưng khác nhau.

\subsection{The Likelihood Problem}

\subsubsection*{Log-likelihood with latent variables}

Với tập dữ liệu gồm $N$ điểm quan sát độc lập, hàm
likelihood đo lường mức độ mà các tham số của mô hình
giải thích được dữ liệu quan sát. Dựa trên giả định các
điểm dữ liệu độc lập, likelihood của toàn bộ tập dữ liệu
được tính bằng tích xác suất mật độ của từng điểm:

\begin{equation}
L(\boldsymbol{\theta})
=
\prod_{n=1}^{N}
p(\bm{x}_n\mid\boldsymbol{\theta})
\end{equation}

Thay công thức phân phối hỗn hợp Gaussian vào likelihood,
ta có:

\begin{equation}
L(\boldsymbol{\theta})
=
\prod_{n=1}^{N}
\left[
\sum_{k=1}^{K}
\pi_k
\mathcal{N}(\bm{x}_n\mid
\bm{\mu}_k,\boldsymbol{\Sigma}_k)
\right]
\end{equation}

Mục tiêu của quá trình ước lượng tham số là tìm ra bộ
tham số làm cho likelihood đạt giá trị lớn nhất:

\begin{equation}
\boldsymbol{\theta}^{*}
=
\arg\max_{\boldsymbol{\theta}}
L(\boldsymbol{\theta})
\end{equation}

Tuy nhiên, việc làm việc trực tiếp với tích của nhiều
mật độ xác suất có thể gây bất tiện về mặt tính toán số.
Vì vậy, người ta thường lấy logarithm của likelihood
để xây dựng hàm log-likelihood:

\begin{align}
\ell(\boldsymbol{\theta})
&= \log L(\boldsymbol{\theta})\\
&= \sum_{n=1}^{N}
\log p(\bm{x}_n\mid\boldsymbol{\theta})\\
&= \sum_{n=1}^{N}
\log\left[
\sum_{k=1}^{K}
\pi_k
\mathcal{N}(\bm{x}_n\mid
\bm{\mu}_k,\boldsymbol{\Sigma}_k)
\right]
\end{align}

Phép logarithm chuyển tích các biểu thức xác suất thành
tổng, giúp việc phân tích và tối ưu hóa thuận tiện hơn.
Trong quá trình triển khai, tính toán log-likelihood
cũng có thể giúp hạn chế hiện tượng tràn số hoặc mất độ
chính xác khi nhân nhiều giá trị rất nhỏ.

\subsubsection*{Intractability of direct solution}

Đối với GMM, biểu thức log-likelihood chứa logarithm
của tổng các thành phần Gaussian. Do đó, việc tối ưu hóa
trực tiếp các tham số có thể trở nên phức tạp.
Phương pháp Expectation-Maximization (EM) thường được
sử dụng để ước lượng các tham số của mô hình thông qua
hai bước lặp: Expectation (E-step) và Maximization
(M-step).

\end{document}